%% kapa: exam=1 round=24 year=2013 period=2 form=A subject=appraisal_law subject_ko=부동산관계법규 seq=17 number=57
%% source: Q-Net 감정평가사 기출문제 (cst00309) artlSeq=5073802 file=extracted/1008/5073802/2013년 제24회 감정평가사 1차시험 기출문제/2교시(A형).hwp
%% conversion: kordoc 4.24.51 -> tools/gen1.py (auto); review=proofread-sonnet; flags=none
\begin{question}{57}
부동산 가격공시 및 감정평가에 관한 법령상의 용어에 관한 설명으로 옳지 않은 것은?
\begin{choices}
  \item 단독주택은 공동주택을 제외한 주택을 말한다.
  \item ‘토지등’에는 어업권과 유가증권이 포함된다.
  \item 지가현황도면이란 당해연도의 산정지가, 전년도의 개별공시지가 및 당해연도의 표준지공시지가가 필지별로 기재된 도면을 말한다.
  \item 적정가격이란 당해 토지 및 주택에 대하여 통상적인 시장에서 정상적인 거래가 이루어지는 경우 성립될 가능성이 가장 높다고 인정되는 가격이다.
  \item 지하주차장 면적을 제외하고 주택으로 쓰이는 1개 동의 연면적이 660제곱미터이고 4개층으로 이루어진 주택은 연립주택이다.
\end{choices}
\end{question}
\answer{5}
